\documentclass[12pt,a4paper]{article}

% --- PAQUETES DE CONFIGURACIÓN BÁSICA ---
\usepackage{babel}
\usepackage[utf8]{inputenc}
\usepackage[T1]{fontenc}
\usepackage[margin=2.2cm]{geometry}
\usepackage{hyperref}
\usepackage{enumitem}
\usepackage{booktabs}
\usepackage{xcolor}
\usepackage{titlesec}
\usepackage{titling}
\usepackage{microtype}
\usepackage{array}
\usepackage{multirow}
\usepackage{listings}
\usepackage{tcolorbox}
\usepackage{float}
\usepackage{lmodern} 

% --- DEFINICIÓN DE COLORES CORPORATIVOS ---
\definecolor{primary}{RGB}{15,118,110}   
\definecolor{secondary}{RGB}{37,99,235} 
\definecolor{muted}{RGB}{107,114,128}     
\definecolor{bgcode}{RGB}{248,250,252}    
\definecolor{frame}{RGB}{226,232,240}     

% --- ESTILO DE SECCIONES ---
\titleformat{\section}{\Large\bfseries\color{primary}}{\thesection}{0.6em}{}
\titleformat{\subsection}{\large\bfseries\color{secondary}}{\thesubsection}{0.6em}{}
\titleformat{\subsubsection}{\bfseries\color{secondary}}{\thesubsubsection}{0.6em}{}

% --- CONFIGURACIÓN DE LISTINGS ---
\lstdefinelanguage{SQLPlus}[]{SQL}{
  morekeywords={IF,ELSE,END,DECLARE,BEGIN,EXCEPTION,WHEN,LOOP,ELSIF,THEN,RAISE,RETURN,ROLE,ACCOUNT,LOCK,UNLOCK,CASE,WHEN,THEN,ELSE,END,CAST,AS,STR_TO_DATE,IFNULL,COALESCE,ALTER,MODIFY,CHANGE,COLUMN,ADD,CONSTRAINT,CHECK,REGEXP},
  sensitive=false
}

\lstset{
  language=SQLPlus,
  basicstyle=\ttfamily\small,
  keywordstyle=\color{secondary}\bfseries,
  commentstyle=\color{muted}\itshape,
  stringstyle=\color{primary},
  columns=fullflexible,
  backgroundcolor=\color{bgcode},
  frame=single,
  rulecolor=\color{frame},
  breaklines=true,
  tabsize=2,
  showstringspaces=false,
  inputencoding = utf8,
  extendedchars = true,
  literate = 
      {á}{{\'a}}1  {é}{{\'e}}1  {í}{{\'i}}1 {ó}{{\'o}}1  {ú}{{\'u}}1
      {Á}{{\'A}}1  {É}{{\'E}}1  {Í}{{\'I}}1 {Ó}{{\'O}}1  {Ú}{{\'U}}1
      {ñ}{{\~n}}1  {Ñ}{{\~N}}1  {¿}{{?`}}1  {¡}{{!`}}1
}

% --- CONFIGURACIÓN DE CAJAS DE TEXTO ---
\tcbset{
  colback=white,
  colframe=primary,
  coltitle=white,
  colbacktitle=primary,
  fonttitle=\bfseries,
  arc=2mm,
  boxrule=0.6pt
}

% --- DATOS DEL DOCUMENTO ---
\title{\textbf{Guía de Teoría: Actualización de Bases de Datos (RA UPDATE)}}
\author{Departamento de Informática}
\date{\today}

\begin{document}

\maketitle

\section*{Introducción}
Esta guía detalla las técnicas y comandos esenciales para la actualización, manipulación y saneamiento de datos en sistemas gestores de bases de datos. 
%
Se abordarán desde los escenarios básicos de inicialización hasta técnicas avanzadas de migración, gestión de colisiones, saneamiento profundo, auditoría y reestructuración del esquema.

\section{PASO 0: El Escenario del Crimen (Script de Inicialización)}
Para seguir esta guía, es necesario crear una base de datos de pruebas que contenga datos ``sucios'' y errores comunes de importación.

\textbf{Instrucciones:} Ejecuta el script adjunto \texttt{00\_fugas\_setup.sql} en tu servidor local para preparar el entorno de trabajo. 
%
Este script genera tablas de producción, tablas de \textit{staging} y registros con inconsistencias deliberadas.

\lstinputlisting[caption={Script de inicialización: \texttt{00\_fugas\_setup.sql}}]{src/00_fugas_setup.sql}

\section{El Entorno de Trabajo y Seguridad}
Para manipular grandes volúmenes de datos, es crucial controlar tanto la terminal como los clientes gráficos, conociendo sus mecanismos de seguridad.

\subsection{Ejecución desde la Terminal y Modo Seguro}
\begin{tcolorbox}[title=Tip de terminal vs GUI]
El saneamiento masivo no debe hacerse ``a mano''. El \textbf{Modo Seguro} es una restricción del \textbf{cliente} (como MySQL Workbench), no del servidor. Por tanto, los scripts ejecutados directamente por terminal (\texttt{mysql < script.sql}) suelen ejecutarse sin estas restricciones, evitando bloqueos por falta de memoria RAM comunes en entornos gráficos al renderizar scripts de gran tamaño.
\end{tcolorbox}

\subsection{MySQL Workbench y la directiva \textit{Safe Updates}}
Workbench lanza el \textbf{Error Code: 1175} si intentas un \texttt{UPDATE} sin una PK en el \texttt{WHERE}. Como ``forenses'', debemos sortearlo:

\lstinputlisting[caption={Gestión de Safe Updates: \texttt{01\_safe\_updates.sql}}]{src/01_safe_updates.sql}

\begin{tcolorbox}[title=Reflexión 0]
¿Qué sentido tiene el modo seguro y por qué viene activado por defecto en los entornos de aprendizaje y clientes gráficos?
\end{tcolorbox}

\section{Integridad Referencial y Restricciones}
Las tablas no son islas; están unidas por claves foráneas (FK) que evalúan reglas cada vez que ejecutamos DML.

\begin{tcolorbox}[title=Recomendación]
Para una comprensión profunda de las FKs, consulta la guía específica de \textbf{Integridad Referencial} disponible en la carpeta \texttt{/home/vicdejuan/BasesDeDatos/Teoria/UT4/01\_On\_delete\_or\_update/}.
\end{tcolorbox}

\subsection{Comportamientos ante UPDATE y DELETE}
\begin{itemize}
    \item \textbf{RESTRICT / NO ACTION:} El escudo principal. Bloquea la sentencia si existen registros dependientes.
    \item \textbf{CASCADE:} El bisturí automático. 
    \begin{itemize}
        \item \texttt{ON DELETE CASCADE}: Borrar un padre fulmina automáticamente a sus hijos (Peligro de pérdida masiva).
        \item \texttt{ON UPDATE CASCADE}: Si el ID principal cambia (ej. corregir un DNI mal metido), el cambio se propaga a todas las tablas dependientes.
    \end{itemize}
    \item \textbf{SET NULL:} Conservación de históricos. Si eliminamos un comercial, no borramos sus ventas (\texttt{CASCADE}), sino que las dejamos ``huérfanas'' pero intactas colocando \texttt{NULL} en su FK.
\end{itemize}

\subsection{Desactivación de Claves Foráneas para Ingesta}
En migraciones, importamos tablas en órdenes aleatorios. La técnica estándar es suspender temporalmente la comprobación:

\lstinputlisting[caption={Control de FKs: \texttt{02\_fk\_checks.sql}}]{src/02_fk_checks.sql}

\begin{tcolorbox}[title=Reflexión 1]
¿Cuál es el problema de eliminar el contenido de las tablas desactivando las claves foráneas (\textit{foreign keys})?
\end{tcolorbox}

\begin{tcolorbox}[colback=red!5!white,colframe=red!75!black,title=\textbf{ADVERTENCIA}]
¡Ojo! Acabas de eliminar todos los datos de la base de datos. Antes de seguir con el resto de ejemplos, vuelve a cargar la parte del script de inserción.
\end{tcolorbox}

\section{Inserción Avanzada y Consolidación}
La inserción fila a fila carece de utilidad en entornos de Big Data. Debemos dominar las inserciones masivas.

\subsection{Motivación: ¿Cuándo necesitamos migrar datos?}
\begin{itemize}
    \item \textbf{Legacy Systems:} Pasar datos de una base antigua a una nueva estructura.
    \item \textbf{Consolidación:} Unificar bases tras una fusión de empresas.
    \item \textbf{ETL de staging:} Procesar datos de Excels externos antes de llevarlos a producción.
\end{itemize}

\subsection{Migración con SELECT y Gestión de Colisiones}
La técnica \texttt{INSERT INTO ... SELECT} permite transformar datos al vuelo.

\lstinputlisting[caption={Migración con SELECT: \texttt{03\_insert\_select.sql}}]{src/03_insert_select.sql}

\begin{tcolorbox}[title=Punto Crítico: Restricciones UNIQUE]
El motor evalúa las restricciones de \texttt{UNIQUE} para cada fila. Si hay una colisión, la sentencia fallará atómicamente a menos que usemos \texttt{IGNORE} o \texttt{ON DUPLICATE KEY UPDATE} (\textit{Upsert}).
\end{tcolorbox}

\lstinputlisting[caption={Estrategia Upsert: \texttt{04\_upsert.sql}}]{src/04_upsert.sql}

\begin{tcolorbox}[title=Reflexión 2]
¿Qué objetivo técnico perseguimos al combinar \texttt{TRIM()}, \texttt{UPPER()} y \texttt{CONCAT()} durante una importación masiva?
\end{tcolorbox}

\begin{tcolorbox}[title=Reflexión 3]
¿Qué ventajas aporta realizar una consulta (\texttt{SELECT}) dentro de una sentencia de inserción (\texttt{INSERT}) en lugar de hacerlo en dos pasos separados?
\end{tcolorbox}

\section{UPDATE: Saneamiento Profundo}
\subsection{Funciones de cadena y Normalización}
Los datos heredados suelen tener espacios excesivos o formatos incorrectos. Usamos funciones nativas:

\lstinputlisting[caption={Normalización de datos: \texttt{05\_phone\_norm.sql}}]{src/05_phone_norm.sql}

\subsection{UPDATE con subconsultas y JOINs}
Para corregir datos basándonos en otras tablas, podemos usar subconsultas o la sintaxis de JOIN múltiple de MySQL.

\lstinputlisting[caption={Update con Subconsultas: \texttt{06\_update\_subquery.sql}}]{src/06_update_subquery.sql}
\lstinputlisting[caption={Update con JOIN: \texttt{07\_update\_join.sql}}]{src/07_update_join.sql}

\subsection{Lógica Condicional Compleja (\texttt{CASE})}
\lstinputlisting[caption={Lógica condicional: \texttt{08\_case\_logic.sql}}]{src/08_case_logic.sql}

\begin{tcolorbox}[title=Reflexión 4]
¿Por qué es preferible usar un \texttt{CASE} dentro de un único \texttt{UPDATE} en lugar de ejecutar varios \texttt{UPDATE} con diferentes cláusulas \texttt{WHERE}?
\end{tcolorbox}

\subsection{El Problema de los \texttt{NULLs}}
Concatenar con \texttt{NULL} devuelve \texttt{NULL}. Usamos \texttt{IFNULL()} o \texttt{COALESCE()} para evitarlo.

\lstinputlisting[caption={Gestión de NULLs: \texttt{09\_null\_handling.sql}}]{src/09_null_handling.sql}

\begin{tcolorbox}[title=Reflexión 5]
Si intentamos concatenar el nombre 'Juan' con un teléfono que es \texttt{NULL}, ¿qué resultado obtenemos y cómo lo evitamos?
\end{tcolorbox}

\subsection{Saneamiento de Fechas y Tipos (\texttt{CAST})}
\lstinputlisting[caption={Saneamiento de tipos: \texttt{10\_date\_saneamiento.sql}}]{src/10_date_saneamiento.sql}

\section{Reestructuración Quirúrgica: Modificación del Esquema (DDL)}
\subsection{Añadir y Eliminar Columnas}
\lstinputlisting[caption={Modificación de columnas: \texttt{11\_modify\_columns.sql}}]{src/11_modify_columns.sql}

\subsection{Modificar Tipos de Datos y Nombres}
\lstinputlisting[caption={Ajuste de tipos de datos: \texttt{12\_modify\_types.sql}}]{src/12_modify_types.sql}

\subsection{Imponer la Integridad: Unicidad y Obligatoriedad}
\lstinputlisting[caption={Añadir Unique y Not Null: \texttt{13\_add\_constraints.sql}}]{src/13_add_constraints.sql}

\begin{tcolorbox}[title=Reflexión 8]
¿Por qué suele fallar un \texttt{ALTER TABLE ... MODIFY ... NOT NULL} en una base de datos que ya contiene datos? ¿Qué paso previo es obligatorio realizar?
\end{tcolorbox}

\subsection{Validación de Reglas de Negocio (\texttt{CHECK})}
\lstinputlisting[caption={Restricciones de validación: \texttt{14\_add\_checks.sql}}]{src/14_add_checks.sql}

\section{Optimización y Rendimiento}
\subsection{SARGability: El peligro de las funciones en el WHERE}
Usar funciones sobre columnas indexadas impide al motor realizar un \textit{Index Seek}.

\lstinputlisting[caption={Rendimiento y SARGability: \texttt{15\_performance\_sarg.sql}}]{src/15_performance_sarg.sql}

\begin{tcolorbox}[title=Reflexión 6]
¿Qué diferencia de rendimiento existe entre \texttt{WHERE UPPER(email) = '...'} y \texttt{WHERE email = '...'} si la columna email está indexada?
\end{tcolorbox}

\section{Eliminación de Datos y Auditoría}
\subsection{DELETE vs TRUNCATE: El Duelo de Gigantes}
\begin{itemize}
    \item \textbf{DELETE:} Operación DML fila a fila. Comprueba restricciones, dispara triggers y escribe cada fila en el log de transacciones. \textbf{Mantiene} el \texttt{AUTO\_INCREMENT}. Es muy lento en tablas masivas.
    \item \textbf{TRUNCATE:} Operación DDL estructural. Destruye y recrea la tabla. Instantánea, no satura el log, pero \textbf{resetea} el \texttt{AUTO\_INCREMENT} y no activa triggers ni admite \texttt{WHERE}.
\end{itemize}

\lstinputlisting[caption={Borrado condicional: \texttt{16\_delete\_date.sql}}]{src/16_delete_date.sql}
\lstinputlisting[caption={Vaciado estructural: \texttt{17\_truncate.sql}}]{src/17_truncate.sql}

\subsection{Verificación Post-Intervención}
\lstinputlisting[caption={Auditoría y Verificación: \texttt{18\_audit\_verify.sql}}]{src/18_audit_verify.sql}

\begin{tcolorbox}[title=Reflexión 7]
¿Por qué es vital usar \texttt{GROUP BY} y \texttt{HAVING} tras un proceso de limpieza de duplicados?
\end{tcolorbox}

\newpage
\section{Respuestas a las Reflexiones e Investigación}

\begin{tcolorbox}[title=Respuesta 0: El Escudo de Seguridad]
Evita borrados o modificaciones accidentales masivas por errores tipográficos en clientes gráficos. \\
\textbf{Pista de investigación:} Busca desastres reales como el de \textit{GitLab} en 2017.
\end{tcolorbox}

\begin{tcolorbox}[title=Respuesta 1: Integridad y Dependencias Circulares]
Al desactivar FKs, el motor permite borrar registros que tienen dependencias activas, lo que puede dejar datos "huérfanos". 
%
Además, es la única forma de vaciar tablas con \textbf{dependencias circulares} (ej. un laboratorio dirigido por un investigador que a su vez pertenece a ese laboratorio). \\
\textbf{Investiga:} Grafos de Dependencia y herramientas como \textit{Flyway}.
\end{tcolorbox}

\begin{tcolorbox}[title=Respuesta 2: ETL y Calidad de Búsqueda]
\textbf{Limpieza:} \texttt{TRIM()} evita errores de duplicados por espacios invisibles. \\
\textbf{Normalización:} \texttt{UPPER()} asegura consistencia en búsquedas posteriores, crítico si el origen de datos no está controlado.
\end{tcolorbox}

\begin{tcolorbox}[title=Respuesta 3: Atomicidad y Eficiencia]
Evita mover datos hacia la aplicación cliente y garantiza que la transformación ocurra de forma atómica y rápida dentro del motor. \\
\textbf{Investiga:} Operaciones Set-based.
\end{tcolorbox}

\begin{tcolorbox}[title=Respuesta 4: Eficiencia de Proceso]
Reduce el tiempo de bloqueo de filas (\textit{Row Locks}) y escanea la tabla una sola vez. \\
\textbf{Pista de investigación:} Investiga los \textbf{Bloqueos de Fila}.
\end{tcolorbox}

\begin{tcolorbox}[title=Respuesta 5: El vacío del NULL]
Obtenemos \texttt{NULL}. Se evita usando \texttt{COALESCE(tel, '')}. \\
\textbf{Pista de investigación:} Busca la \textbf{Lógica Tri-valuada} de SQL.
\end{tcolorbox}

\begin{tcolorbox}[title=Respuesta 6: SARGability]
\texttt{UPPER(email)} obliga a un \textit{Full Table Scan}, mientras que la comparación directa permite un \textit{Index Seek}. \\
\textbf{Pista de investigación:} Busca el término \textbf{SARGable queries}.
\end{tcolorbox}

\begin{tcolorbox}[title=Respuesta 7: Validación de Calidad]
Permite confirmar que no han quedado registros duplicados. Es el control de calidad final. \\
\textbf{Pista de investigación:} Investiga técnicas de \textbf{Data Profiling}.
\end{tcolorbox}

\begin{tcolorbox}[title=Respuesta 8: El cerrojo de la obligatoriedad]
Falla si existen filas con valores \texttt{NULL}. El paso previo es sanear esos nulos mediante un \texttt{UPDATE} antes de aplicar la restricción.
\end{tcolorbox}

\end{document}
